% =====================================================================
\chapter{Independencia lineal, bases y rango}
\label{tema:1}
% =====================================================================

Este tema introduce los conceptos que permiten medir cuántas dimensiones utiliza realmente una matriz: independencia lineal, base, dimensión, rango y núcleo. La presentación sigue los §3.1, §3.4 y~§3.5 de Strang~\cite{strang-ila}, las clases~9 a~11 del curso 18.06~\cite{mit1806} y los §2.5 a~§2.7 de~\cite{mml}. Se trabaja siempre en $\R^{n}$, aunque las definiciones valen en cualquier espacio vectorial real.

\section{Combinaciones lineales y subespacios}

\begin{definicion}[Subespacio]
Un subconjunto $W\subseteq\R^{n}$ es un \emph{subespacio vectorial} si contiene el vector nulo y es cerrado para la suma y el producto por escalares: si $u,v\in W$ y $c\in\R$, entonces $u+v\in W$ y $cu\in W$.
\end{definicion}

\begin{definicion}[Subespacio generado]
Dados $v_1,\dots,v_k\in\R^{n}$, una \emph{combinación lineal} de ellos es un vector $c_1v_1+\dots+c_kv_k$ con $c_i\in\R$. El conjunto de todas las combinaciones lineales es un subespacio, que se denota $\gen\{v_1,\dots,v_k\}$ y se llama \emph{subespacio generado} por los $v_i$.
\end{definicion}

\begin{definicion}[Independencia lineal]\label{def:indep}
Los vectores $v_1,\dots,v_k$ son \emph{linealmente independientes} si la única combinación lineal nula es la trivial:
\[
c_1v_1+\dots+c_kv_k=0\ \Longrightarrow\ c_1=\dots=c_k=0.
\]
En caso contrario se dicen \emph{linealmente dependientes}.
\end{definicion}

Si $V=[v_1\mid\cdots\mid v_k]\in\R^{n\times k}$ es la matriz cuyas columnas son los vectores, la condición equivale a que el sistema homogéneo $Vc=0$ tenga solo la solución $c=0$.

\begin{proposicion}\label{prop:dependencia}
Los vectores $v_1,\dots,v_k$ ($k\ge 2$) son linealmente dependientes si y solo si alguno de ellos es combinación lineal de los demás.
\end{proposicion}

\section{Bases, dimensión y coordenadas}

\begin{definicion}[Base]
Una \emph{base} de un subespacio $W$ es una familia de vectores de $W$ linealmente independiente que genera $W$.
\end{definicion}

\begin{proposicion}[Coordenadas]\label{prop:coordenadas}
Si $\{u_1,\dots,u_k\}$ es una base de $W$, todo $w\in W$ se escribe de forma única como $w=c_1u_1+\dots+c_ku_k$. Los números $c_1,\dots,c_k$ son las \emph{coordenadas} de $w$ en esa base.
\end{proposicion}

\begin{lema}\label{lem:steinitz}
Si $W=\gen\{v_1,\dots,v_m\}$ y $w_1,\dots,w_p\in W$ con $p>m$, entonces $w_1,\dots,w_p$ son linealmente dependientes.
\end{lema}

\begin{teorema}[Dimensión]\label{teo:dimension}
Todas las bases de un subespacio $W$ tienen el mismo número de elementos. Ese número se llama \emph{dimensión} de $W$ y se denota $\dim W$. En particular, $\dim\R^{n}=n$.
\end{teorema}

\addtocounter{definicion}{1}%

\begin{observacion}
Las coordenadas de un embedding dependen de la base en la que se expresan. Un cambio de base invertible del espacio de activaciones de un modelo, compensado en las capas siguientes, no altera sus predicciones; por eso las coordenadas individuales de un embedding no tienen, en general, un significado propio. Esta observación es el punto de partida de los métodos de interpretabilidad del módulo~23, que buscan direcciones del espacio y no coordenadas.
\end{observacion}

\section{Espacio columna, núcleo y rango}

\begin{definicion}[Espacio columna y núcleo]
Sea $A\in\R^{m\times n}$. El \emph{espacio columna} (o imagen) de $A$ es $\Col(A)=\{Ax : x\in\R^{n}\}\subseteq\R^{m}$, el subespacio generado por sus columnas. El \emph{núcleo} de $A$ es $\Nuc(A)=\{x\in\R^{n}: Ax=0\}\subseteq\R^{n}$.
\end{definicion}

\begin{definicion}[Rango]
El \emph{rango} de $A$ es $\rango(A)=\dim\Col(A)$, el número máximo de columnas linealmente independientes de~$A$.
\end{definicion}

\begin{teorema}[Rango de filas igual a rango de columnas]\label{teo:filas-columnas}
Para toda $A\in\R^{m\times n}$, $\rango(A)=\rango(A\T)$: el número máximo de filas independientes coincide con el de columnas independientes.
\end{teorema}

La factorización $A=CR$ de la demostración es la que Strang usa como punto de partida en~\cite[§1.4]{strang-ila}.

\begin{teorema}[Rango y nulidad]\label{teo:rango-nulidad}
Para toda $A\in\R^{m\times n}$,
\[
\rango(A)+\dim\Nuc(A)=n.
\]
\end{teorema}

\addtocounter{definicion}{1}%

\section{Matrices de rango 1 y rango de un producto}

\begin{proposicion}\label{prop:rango1}
Si $u\in\R^{m}$ y $v\in\R^{n}$ son no nulos, la matriz $uv\T\in\R^{m\times n}$ tiene rango~1. Recíprocamente, toda matriz de rango~1 es de esta forma.
\end{proposicion}

\begin{proposicion}[Rango de un producto]\label{prop:rango-producto}
Si $B\in\R^{m\times r}$ y $A\in\R^{r\times n}$, entonces $\rango(BA)\le\min\{\rango(A),\rango(B)\}\le r$.
\end{proposicion}

La demostración se propone en el \cref{eje:1-producto}. De la factorización $A=CR$ del \cref{teo:filas-columnas} se deduce además que toda matriz de rango $r$ es suma de $r$ matrices de rango~1, $A=\sum_{i=1}^{r}c_i\,\rho_i\T$, donde $c_i$ es la columna $i$ de $C$ y $\rho_i\T$ la fila $i$ de $R$. El tema~4 obtiene una descomposición de este tipo con piezas ortogonales y ordenadas por importancia.

\begin{observacion}[Matrices de bajo rango en un LLM]\label{obs:lora-rango}
El ajuste fino LoRA~\cite{hu-lora}, que se estudia en el módulo~20, mantiene congelada una matriz de pesos $W_0\in\R^{d_{\text{out}}\times d_{\text{in}}}$ y aprende una corrección $\Delta W=BA$ con $B\in\R^{d_{\text{out}}\times r}$ y $A\in\R^{r\times d_{\text{in}}}$. Por la \cref{prop:rango-producto}, $\rango(\Delta W)\le r$, y el número de parámetros entrenables es $r(d_{\text{in}}+d_{\text{out}})$ en lugar de $d_{\text{in}}d_{\text{out}}$. Con $d_{\text{in}}=d_{\text{out}}=4096$ y $r=8$ se entrenan $65\,536$ parámetros, el $0{,}39\,\%$ de los $16\,777\,216$ de la matriz completa.
\end{observacion}

